\begin{claim}{C.86}{Proved}{The exact spiral centralizer comparison and its based sphere fillings}
Retain the actual radial coordinate $j$ of C.83--C.85
and the actual spiral coordinate $p$ of C.72:
\[
 j(z)=\frac{z}{1+|z|},\qquad
 p(z)=\frac{z}{1+|z|}e^{i\log(1+|z|)},\qquad
 \mathcal R(z)=iz.
                                                               \tag{C.86.1}
\]
Let $\mathscr C_{\mathrm{rad}}$ be the relative
closed-disk centralizer of $jT_1j^{-1}$, and let
$\mathscr C_{\mathrm{sp}}$ be the relative
closed-disk centralizer of the original C.72
operator $pT_i p^{-1}$. All these operators have
their specified identity extensions on the circle.
There is an actual topological group isomorphism
\[
 \begin{gathered}
 \Theta:\mathscr C_{\mathrm{rad}}
                \longrightarrow\mathscr C_{\mathrm{sp}},\\
 \Theta(h)|_{\mathbb D}
       =p\mathcal Rj^{-1}h j\mathcal R^{-1}p^{-1},
       \qquad \Theta(h)|_{S^1}=\id.
 \end{gathered}                                                   \tag{C.86.2}
\]
Its inverse is given by the inverse interior
conjugation and the same identity boundary
extension. The proof establishes both topology
maps in the original map-and-inverse topology.
In particular,
\[
             \pi_n(\mathscr C_{\mathrm{sp}},1)=0
                          \qquad(n\geq1).                        \tag{C.86.3}
\]
The interior coordinate
$p\mathcal Rj^{-1}$ in (C.86.2) has no
continuous extension at the boundary point
$1$. The group comparison and the stated
fillings are established by the exact
sublinearity and topology calculations below.
\end{claim}

\begin{proof}
\emph{A global spiral displacement estimate.}
For all $z,v\in\mathbb C$ one has
\[
              |p(z+v)-p(z)|
                  \leq\frac{4|v|}{1+|z|}.                       \tag{C.86.4}
\]
Put $r=|z|$, $r'=|z+v|$, and $M=|v|$.
For $M\leq(1+r)/2$, the reverse triangle
inequality gives
$\min\{1+r,1+r'\}\geq(1+r)/2$.
The mean value theorem for the logarithm
and (C.83.3) therefore give
\[
 \begin{split}
 |\log(1+r')-\log(1+r)|&\leq\frac{2M}{1+r},\\
 |p(z+v)-p(z)|
  &\leq |j(z+v)-j(z)|
       +|j(z)|\,|e^{i\log(1+r')}-e^{i\log(1+r)}|\\
  &\leq\frac{4M}{1+r}.
 \end{split}                                                     \tag{C.86.5}
\]
For $M>(1+r)/2$, both spiral values lie
in the unit disk, so their distance
is at most $2<4M/(1+r)$. This proves
the global estimate, including at zero.

We first compare the horizontal operator
$pT_1p^{-1}$ with $jT_1j^{-1}$.
The final rotation will restore the
literal vertical label of C.72.

\emph{Necessity, uniformly on compact actual families.}
Let $K$ be a compact family of original
closed-disk homeomorphisms fixing the
circle and commuting with $pT_1p^{-1}$.
Write their actual plane conjugates as
$F=p^{-1}hp$.
Each $F$ commutes with the unit horizontal
translation. Compactness and boundary
fixation give the exact uniform limit
\[
 \sup_{h\in K,\ |z|\geq R}|pF(z)-p(z)|
                              \longrightarrow0.
                                                               \tag{C.86.6}
\]
Indeed, $|p(z)|\to1$. Any failing
sequence has a subsequence on which
$p(z)$ approaches one circle point
and $h$ converges uniformly to a map
fixing that point. Its displayed
difference then tends to zero.

Take arbitrary $F_i$ from these conjugates
and $(x_i,y_i)$ with $0\leq x_i\leq1$
and $|y_i|\to\infty$. Put
\[
 U_i=\frac{X_{F_i}(x_i,y_i)-x_i}{y_i},\qquad
 V_i=\frac{Y_{F_i}(x_i,y_i)}{y_i}.
                                                               \tag{C.86.7}
\]
For any fixed real $c$ choose the actual
integer $m_i(c)$ so that
\[
 x_i+m_i(c)=cy_i+e_i(c),\qquad |e_i(c)|\leq\tfrac12.
                                                               \tag{C.86.8}
\]
The input and output at this integer
translate are exactly
\[
 \begin{gathered}
 z_i(c)=y_i(c+i)+e_i(c),\\
 F_i(z_i(c))=y_i(c+U_i+iV_i)+e_i(c).
 \end{gathered}                                                   \tag{C.86.9}
\]
Here $i$ in the complex expressions is
the imaginary unit; subscripts retain
the sequence index.

First $(U_i,V_i)$ is bounded. Suppose a
subsequence has norm tending to infinity.
The two outputs for $c=0$ and $c=1$
differ by the real integer
$m_i(1)-m_i(0)$, of absolute value at
most $|y_i|+1$. The norm of the output
at zero, divided by $|y_i|$, tends to
infinity. Applying (C.86.4) to these
two outputs shows that their spiral
images have distance tending to zero.
In contrast, the distance of the two
spiral input images tends to a positive
number. To see its exact value, define
\[
       \Psi(w)=\frac{w}{|w|}e^{i\log|w|}
                         \quad(w\ne0).
                                                               \tag{C.86.10}
\]
After removal of the common unit factor
$\operatorname{sign}(y_i)e^{i\log|y_i|}$,
the two input phases tend to $\Psi(i)$
and $\Psi(1+i)$. Their angle difference
is
$\delta=\pi/4-(\log2)/2$, with
$0<\delta<\pi/4$; the first inequality
follows from $\log2<1<\pi/2$.
Their distance is $2\sin(\delta/2)>0$.
Formula (C.86.6) would force the input
and output distances to have the same
limit, a contradiction. This proves
boundedness without presupposing a
normalized nonzero output limit.

Pass to a subsequence with
$(U_i,V_i)\to(U,V)$. For every real
$c$ with $c+U+iV\ne0$, formula (C.86.9)
and removal of the same common unit
factor give, by (C.86.6),
\[
                     \Psi(c+U+iV)=\Psi(c+i).
                                                               \tag{C.86.11}
\]
This is equality of circle-valued
functions. No equality of real
arguments, or choice of an integer
multiple of $2\pi$, has been imposed.
On every open interval where the
left argument is nonzero, differentiate
the circle-valued identity and divide
by its nonzero value. The exact
logarithmic derivative is
\[
 \frac{d}{dc}\Psi(c+U+iV)\big/\Psi(c+U+iV)
       =i\frac{c+U-V}{(c+U)^2+V^2}.
                                                               \tag{C.86.12}
\]
Thus
\[
 (c+U-V)(c^2+1)
       =(c-1)((c+U)^2+V^2).                                    \tag{C.86.13}
\]
The identity holds as a polynomial
because it holds on an open interval.
This includes the possible limit
$V=0$: only the single value $c=-U$
was omitted, and no conclusion at
that zero was needed.
Comparing the coefficient of $c^2$
gives $U+V=1$; comparing constants
then gives
$U-V+U^2+V^2=2U^2=0$.
Consequently $U=0$ and $V=1$.

Every convergent subsequence of the
bounded pair has this same limit.
Hence
\[
 \frac{|F_i(x_i,y_i)-(x_i,y_i)|}{|y_i|}
           =\sqrt{U_i^2+(V_i-1)^2}\longrightarrow0.
                                                               \tag{C.86.14}
\]
Since the original sequences were
arbitrary, this proves precisely the
uniform compact-family sublinearity
condition (C.84.1) for the spiral
relative centralizer. Apply the same
argument to the compact inverse family
to retain its inverse condition.

\emph{Sufficiency and the two inverse topology maps.}
Conversely, let $F\in\mathscr P$ as
defined in C.84. Periodicity and its
sublinearity give
$|F(z)-z|/(1+|z|)\to0$ at infinity,
as proved in (C.84.6). Formula (C.86.4)
shows that $pFp^{-1}$ extends by the
identity to the original circle.
The extension is a continuous bijection
of the compact disk and thus a
homeomorphism. It commutes with the
specified horizontal spiral operator;
its inverse is the extension of
$pF^{-1}p^{-1}$. We have proved an
exact group bijection
\[
 \begin{gathered}
 I_p:\mathscr P\longrightarrow
        C_{\mathcal D_\partial}(pT_1p^{-1}),\\
 I_p(F)|_{\mathbb D}=pFp^{-1},\qquad
 I_p(F)|_{S^1}=\id,
 \end{gathered}                                                   \tag{C.86.15}
\]
whose inverse is restriction followed
by $p^{-1}(\cdot)p$.

For its forward continuity fix $H\in\mathscr P$.
The finite constant $c_H$ of (C.84.7)
and (C.86.4), applied at $H(z)$, give
\[
 \|I_p(F)-I_p(H)\|_\infty
          \leq4c_H\|F-H\|_*.
                                                               \tag{C.86.16}
\]
The estimate uses the exact periodic
displacement and the original source
height in the definition of $\|\cdot\|_*$.
Applying it to $F^{-1},H^{-1}$ proves
continuity for the original inverses
as well.

For the inverse topology map, a
convergent sequence of actual disk
maps together with its limit is
compact. The necessity proof above
gives uniform sublinear tails for
the corresponding plane maps and
their inverses. On each compact
period rectangle, conjugation by
$p,p^{-1}$ gives uniform convergence:
the disk images stay in one compact
interior set, where $p^{-1}$ is
uniformly continuous. Combining that
convergence with the common tail
bounds gives convergence of the
weighted forward differences and
of the weighted inverse differences.
Thus the inverse of (C.86.15) is
continuous in $\delta_*$.
These metrizable-space arguments prove
the asserted topological group
isomorphism with the actual topology.

\emph{The original vertical label, the nonextendible coordinate, and the filling.}
Both $j$ and $p$ commute with the
unit complex rotation $\mathcal R$.
Its closed-disk extension $\overline{\mathcal R}(w)=iw$
conjugates the horizontal spiral
operator to the original vertical
operator of C.72. The exact C.84
isomorphism $I_j$ and (C.86.15)
therefore give
\[
 \Theta(h)=\overline{\mathcal R}\,
               I_p(I_j^{-1}(h))\,
                    \overline{\mathcal R}^{-1}.
                                                               \tag{C.86.17}
\]
This is (C.86.2) on the interior and
the identity on the boundary.
Its inverse is
$I_j I_p^{-1}(\overline{\mathcal R}^{-1}(\cdot)
\overline{\mathcal R})$.
The two inverse identities and
both continuous topology maps follow
from the specified component maps;
every original integer operator
label is preserved by the displayed
rotation and conjugations.

For completeness, the interior
coordinate itself does not extend
at the indicated boundary point.
Put $r_k=e^{2\pi k}-1$ and
$s_k=e^{(2k+1)\pi}-1$.
Both $j(r_k)$ and $j(s_k)$ tend to
$1$, whereas
\[
 \begin{gathered}
 p\mathcal Rj^{-1}(j(r_k))
        =i\frac{r_k}{1+r_k}\longrightarrow i,\\
 p\mathcal Rj^{-1}(j(s_k))
        =-i\frac{s_k}{1+s_k}\longrightarrow-i.
 \end{gathered}                                                   \tag{C.86.18}
\]
Thus the group isomorphism above
is proved by (C.86.6)--(C.86.17),
with its own actual boundary
continuity, rather than by extending
that coordinate.

Given a based sphere family in
$\mathscr C_{\mathrm{sp}}$, apply
the continuous inverse $\Theta^{-1}$,
then the exact C.85 filling, and
finally $\Theta$ at every homotopy
parameter. This is a continuous
based filling in the original
spiral relative centralizer, for
both maps and inverses, and proves
(C.86.3). It covers the entire
relative centralizer of C.72,
including families beyond its
explicitly contracted example.
The general discrete-operator case
and general C.12 remain open.
\end{proof}
